\documentclass[12pt]{article}
\usepackage[utf8]{inputenc}
\usepackage[spanish]{babel}
\usepackage{graphicx}
\usepackage{hyperref}
\usepackage{geometry}
\usepackage{fancyhdr}
\usepackage{float}
\usepackage{verbatim}

% Configuración de página
\geometry{a4paper, margin=2.5cm}
\setlength{\parindent}{0pt}
\setlength{\parskip}{1em}

% Configuración de encabezados
\pagestyle{fancy}
\fancyhf{}
\rhead{Práctica Git - DAWD}
\lhead{Sistema de Gestión de Notas}
\cfoot{\thepage}

\begin{document}

\begin{titlepage}
    % Logos en las esquinas superiores
    \begin{center}
        \begin{tabular}{@{}p{0.3\textwidth}@{}p{0.5\textwidth}@{}p{0.3\textwidth}@{}}
            \includegraphics[width=0.2\textwidth]{logo_unam.png} & 
            \hfill &
            \includegraphics[width=0.15\textwidth]{logo_iimas.png}
        \end{tabular}
    \end{center}
    
    \vspace*{1cm}
    
    % Título principal
    \begin{center}
        {\Huge \textbf{Universidad Nacional Autónoma de México}}\\[0.5cm]
        {\Large \textbf{Instituto de Investigaciones en Matemáticas Aplicadas y en Sistemas}}\\[2cm]
        
        {\huge \textbf{Práctica 6}}\\[0.5cm]
        {\LARGE Control de Versiones con Git y GitHub}\\[1.5cm]
        
        {\Large \textbf{Alumno:}}\\[0.3cm]
        {\Large Erick José Fabián Sandoval}\\[1cm]
        
        {\Large \textbf{Asignatura:}}\\[0.3cm]
        {\Large Desarrollo de Aplicaciones Web D.}\\[1cm]
        
        {\Large \textbf{Profesor:}}\\[0.3cm]
        {\Large Israel Sandoval Grajeda}\\[1.5cm]

        {\Large  \url{https://github.com/erick0x/P6_DAWD}} \\[1cm]
        
        \vfill
        {\large \today}
    \end{center}
\end{titlepage}

\tableofcontents
\newpage

\section{Introducción}
Este informe documenta la práctica de pruebas unitarias realizadas para una aplicación Angular y un módulo Python. El objetivo es demostrar el desarrollo de casos de prueba que validan la conexión a un servicio web en Angular y la conexión a una base de datos en Python.

\subsection{Objetivos}
\begin{itemize}
    \item Desarrollar una prueba unitaria en Angular para verificar la conexión a un web service.
    \item Desarrollar una prueba unitaria en Python para validar la conexión a una base de datos.
    \item Mantener la portada original del documento.
\end{itemize}

\section{Metodología}
La práctica se realizó en dos etapas principales:
\begin{enumerate}
    \item Implementación del servicio y la prueba unitaria en Angular.
    \item Implementación del módulo de conexión y la prueba unitaria en Python.
\end{enumerate}
En ambos casos se utilizó un archivo de prueba separado del código fuente, tal como se solicita en la especificación de la práctica.

\section{Prueba unitaria en Angular}
La aplicación Angular creada contiene un servicio que prueba la conexión a un web service público. El componente principal invoca este servicio y actualiza un estado interno que indica si la conexión fue exitosa.

\subsection{Código de la prueba}
La prueba se encuentra en el archivo \texttt{src/app/app.spec.ts}. Se utilizan los siguientes elementos:
\begin{itemize}
    \item \texttt{TestBed} para configurar el entorno de pruebas.
    \item \texttt{HttpClientTestingModule} para simular las peticiones HTTP.
    \item \texttt{HttpTestingController} para controlar las respuestas mock.
\end{itemize}

\subsection{Casos evaluados}
\begin{itemize}
    \item \textbf{Conexión exitosa}: cuando el servicio responde correctamente, el estado se actualiza a \texttt{online}.
    \item \textbf{Conexión fallida}: cuando ocurre un error de red, el estado se actualiza a \texttt{offline}.
\end{itemize}

\subsection{Resultado de ejecución}
La prueba Angular se ejecutó con el comando:
\begin{verbatim}
cd angular-app
npx ng test --watch=false
\end{verbatim}

\noindent La salida mostrada fue:
\begin{verbatim}
✓  angular-app  src/app/app.spec.ts (3 tests)
   ✓ App (3)
     ✓ should create the app
     ✓ should set status online when web service responds successfully
     ✓ should set status offline when web service returns an error
\end{verbatim}

\section{Prueba unitaria en Python}
El módulo Python implementa una clase de conexión a una base de datos SQLite. La prueba unitaria se realizó en un archivo separado para garantizar la limpieza entre el código fuente y las pruebas.

\subsection{Código de la prueba}
La prueba se encuentra en el archivo \texttt{python-app/test_database.py} y utiliza el módulo \texttt{unittest}. Las operaciones cubiertas incluyen:
\begin{itemize}
    \item Conexión a la base de datos.
    \item Desconexión.
    \item Creación de tabla.
    \item Inserción de datos.
    \item Consulta de datos.
    \item Manejo de casos sin conexión.
\end{itemize}

\subsection{Resultado de ejecución}
La prueba Python se ejecutó con el comando:
\begin{verbatim}
cd python-app
python -m unittest test_database.py -v
\end{verbatim}

\noindent La salida mostrada fue:
\begin{verbatim}
test_connect_successful (test_database.TestDatabaseConnection) ... ok
test_create_table (test_database.TestDatabaseConnection) ... ok
test_disconnect_successful (test_database.TestDatabaseConnection) ... ok
test_insert_data (test_database.TestDatabaseConnection) ... ok
test_insert_without_connection_returns_false (test_database.TestDatabaseConnection) ... ok
test_query_data (test_database.TestDatabaseConnection) ... ok
test_query_without_connection_returns_none (test_database.TestDatabaseConnection) ... ok
\end{verbatim}

\section{Resultados generales}
Ambas pruebas unitarias se ejecutaron correctamente y cumplieron el propósito de validar las conexiones en cada entorno:
\begin{itemize}
    \item La prueba Angular verifica la comunicación con un servicio web externo mediante un mock HTTP.
    \item La prueba Python verifica la conexión a una base de datos SQLite y el comportamiento ante operaciones básicas.
\end{itemize}

\section{Conclusiones}
La práctica demuestra que es posible implementar pruebas unitarias en diferentes lenguajes y plataformas manteniendo una estructura de código clara. En Angular se utilizó un servicio con una prueba separada en el archivo \texttt{app.spec.ts}; en Python se empleó un archivo de prueba independiente para validar la conexión y operaciones a la base de datos.

\section{Anexos}
\begin{itemize}
    \item Archivo de prueba Angular: \texttt{angular-app/src/app/app.spec.ts}
    \item Servicio Angular: \texttt{angular-app/src/app/web.service.ts}
    \item Código Python de conexión: \texttt{python-app/database.py}
    \item Prueba Python: \texttt{python-app/test_database.py}
\end{itemize}


\subsection{Resultados Obtenidos}
\begin{itemize}
    \item Se implementaron exitosamente las 5 operaciones fundamentales de Git
    \item Se desarrolló una aplicación web funcional con HTML, CSS y JavaScript
    \item Se mantuvo un historial de versiones organizado con 4 commits
    \item Se demostró el flujo de trabajo con ramas (branching strategy)
    \item El repositorio remoto contiene el código completo y documentación
\end{itemize}

\subsection{Comandos Git Utilizados}
\begin{table}[H]
    \centering
    \begin{tabular}{|l|l|}
        \hline
        \textbf{Comando} & \textbf{Propósito} \\
        \hline
        git init & Inicializar repositorio local \\
        git add & Agregar archivos al staging \\
        git commit & Crear commit con cambios \\
        git checkout -b & Crear y cambiar a nueva rama \\
        git merge & Fusionar ramas \\
        git branch -d & Eliminar rama \\
        git push & Subir cambios al remoto \\
        git log --oneline --graph & Ver historial gráfico \\
        git remote -v & Ver repositorios remotos \\
        \hline
    \end{tabular}
    \caption{Comandos Git utilizados en la práctica}
\end{table}

\subsection{Referencias}
\begin{itemize}
    \item Repositorio del proyecto: \url{https://github.com/erick0x/P6_DAWD}
    \item Documentación oficial de Git: \url{https://git-scm.com/doc}
\end{itemize}

\end{document}
